\documentclass[11pt]{article}
\usepackage[T1]{fontenc}
\usepackage{lmodern,amsmath,amssymb,amsthm,mathtools,booktabs,geometry,microtype}
\usepackage[hidelinks]{hyperref}
\hypersetup{pdftitle={The iterated-logarithm law for fixed-cardinality p-adic torsion specialization},pdfauthor={Research note for the ProveIt GRH program}}
\setlength{\emergencystretch}{2em}
\geometry{margin=1.05in}
\setlength{\parskip}{4pt}
\setlength{\parindent}{0pt}
\newtheorem{theorem}{Theorem}[section]
\newtheorem{lemma}[theorem]{Lemma}
\newtheorem{corollary}[theorem]{Corollary}
\newtheorem{proposition}[theorem]{Proposition}
\theoremstyle{remark}\newtheorem{remark}[theorem]{Remark}
\newcommand{\Cp}{\mathbb C_p}
\newcommand{\OO}{\mathcal O}
\newcommand{\Z}{\mathbb Z}
\newcommand{\Q}{\mathbb Q}
\newcommand{\G}{G_\rho}
\newcommand{\lp}{\log_p}
\newcommand{\pos}[1]{\left[#1\right]_+}
\newcommand{\code}[1]{\texttt{#1}}
\title{The iterated-logarithm law for fixed-cardinality\\ $p$-adic torsion specialization}
\author{Research note for the ProveIt GRH program}
\date{27 November 2026}
\begin{document}
\maketitle
\begin{abstract}
The torsion-selection step of the vendored Dirichlet GRH proof controls the valuation increase of a finite packet of polynomials at a common tuple of prime-power roots of unity. The merged ProveIt article determines exact obstruction costs for one polynomial and for packets of unrestricted size, but asks for the intervening fixed-cardinality problem. We prove a sharp iterated-logarithmic asymptotic for every fixed packet cardinality at zero tolerance. With budget $B$, starting shell $h_0$, and at most $n$ polynomials, the least guaranteed order exponent $H_n(B)$ satisfies
\[
 \log_p^{\circ(n-1)} H_n(B)=p^{h_0-1} B+O_{p,h_0,n}(\log(2+B)).
\]
The lower error is bounded independently of $B$; for $n=1$ the answer is exact. The upper bound comes from a delayed product potential, which avoids losing a factor $p/(p-1)$ at each stage of the argument. The theorem is degree-free, dimension-free, and valid for integral power series. We also give an exact finite integer-state algorithm for the obstruction cost at every nonnegative tolerance, a stabilization result at tolerance zero, and a Dirichlet-GRH application to simultaneous character specialization at a bounded prime conductor. The local theorem itself is unconditional. Neither the upstream Lean proof nor these new results are independently kernel-checked in this note.
\end{abstract}

\section{The problem extracted from the source}
The pinned vendored Lean development \cite{vendor,leanreadme} is described as the readable edition of Levent Alp\"oge's proof, rewritten with Claude, at upstream commit\newline \code{72fead102f663d5c7a407964a13dc0cf9d977207}. Its public interface \cite{statements,main} gives Mathlib's Riemann hypothesis and nonvanishing of Dirichlet $L$-functions to the right of the critical line, with the principal pole excluded. The README reports standard Lean axioms only and independent checker runs; this note uses those statements as supplied, not as a freshly checked build. In particular we do not infer GRH for arbitrary Artin or automorphic $L$-functions.

The relevant part of the proof is not the last analytic zero test. In \code{MainInequality.lean}, the upper and lower estimates for a corrected logarithmic covolume cancel their common main term; \code{OneSidedBound.lean} converts the remaining one-sided prime estimate into Dirichlet nonvanishing \cite{inequality,onesided}. The lower covolume estimate uses favorable twists, controlled via Selmer groups and valuations of ratios of Hecke values. The file \code{RootsOfUnity.lean} supplies degree-independent local specialization lemmas:
\begin{itemize}
\item \code{exists\_rootOfUnity\_eval\_le\_gaussVal\_add}, a one-variable result based on root counting;
\item \code{exists\_rootsOfUnity\_ellVal\_eval\_le\_add}, a common tuple for multiple polynomials and multiple variables.
\end{itemize}
The file \code{TwistingTheorem.lean} describes the further arithmetic constraints on a common tame twist \cite{roots,twisting}. The local selector has a positive tolerance parameter. The workbook's chapter on roots of unity for many polynomials also explains the harmonic root-count argument \cite{workbook}.

The merged article \cite{article,artcost,arttors} states two different minimax answers: for a single polynomial, the exact zero-tolerance cost of blocking consecutive shells is linear in the number of shells; for unrestricted packets the cost stays bounded at zero tolerance and is asymptotically harmonic for positive tolerance. In \code{cost.tex}, the question \code{cost:q:cardinality} explicitly asks for the least blocking cost with a prescribed number of members. The question \code{tors:q:cardinality} asks whether the bounds for fixed families with separate tolerances can be improved. We address the first question asymptotically at zero tolerance, and supply an exact finite algorithm for both questions' common obstruction model. This is a transfer of a mathematical ingredient of the proof, not a new proof of GRH or a claim to improve the arithmetic constants of the original argument without further work.

\section{Definitions and main results}
Fix a prime $p$, the completed algebraic closure $\Cp$ of $\Q_p$, and its additive valuation $v$ normalized by $v(p)=1$, with $v(0)=+\infty$. Write $\OO=\{x:v(x)\ge0\}$. For $h\ge1$ put
\[
 e_h=(p-1)p^{h-1},\quad r_h=e_h^{-1},\quad
 Z_h=\{\zeta-1: \zeta\text{ has exact order }p^h\}.
\]
Then $\#Z_h=e_h$ and $v(z)=r_h$ on $Z_h$. Fix $h_0\ge1$ and $\rho\ge r_{h_0}$. Define $s_0=p^{1-h_0}$. For a nonzero polynomial $g=\sum_j c_jT^j\in\OO[T]$ let $G_r(g)=\min_j(v(c_j)+jr)$ for $r\ge0$.

A packet of at most $n$ members consists of nonzero integral polynomials $g_i$ and finite real thresholds $t_i\ge G_\rho(g_i)$. Its budget is $\sum_i t_i$ and its total positive loss at $z$ is
\[
 \mathcal L(z)=\sum_i\pos{v(g_i(z))-t_i},\qquad [x]_+=\max(x,0).
\]
It blocks $[h_0,H]$ at tolerance $\delta\ge0$ if $\mathcal L(z)>\delta$ at every point in every shell. At zero tolerance a point is good precisely when \emph{every} individual constraint holds. Let $C_n(H,\delta)$ denote the infimum budget of a blocker, with $p,h_0,\rho$ suppressed. Let $H_n(B,\delta)$ be the least integer $H\ge h_0$ such that every packet of budget at most $B$ has some $z\in\bigcup_{h=h_0}^H Z_h$ with $\mathcal L(z)\le\delta$. Write $H_n(B)=H_n(B,0)$. Unused members may be padded with the constant $1$ at threshold zero. Superscripts $\log_p^{\circ j}$ mean $j$ compositions, with the zeroth composition the identity.

\begin{theorem}[Sharp fixed-cardinality tower law]\label{thm:main}
For fixed $p,h_0,n$ and every $B\ge0$, $H_n(B)$ is finite, uniformly in $\rho\ge r_{h_0}$. For $n=1$,
\[
 H_1(B)=h_0+\lfloor B/s_0\rfloor.
\]
For fixed $n\ge2$, as $B\to\infty$,
\begin{equation}\label{eq:tower}
 \frac{B}{s_0}-O_{p,h_0,n}(1)\ \le\ \log_p^{\circ(n-1)}H_n(B)
 \ \le\ \frac{B}{s_0}+O_{p,h_0,n}(\log(2+B)).
\end{equation}
Equivalently, with $X_H=\log_p^{\circ(n-1)}(H+1)$, as $H\to\infty$,
\begin{equation}\label{eq:inverse}
 s_0X_H-O_{p,h_0,n}(\log(2+X_H))\le C_n(H,0)\le s_0X_H+O_{p,h_0,n}(1).
\end{equation}
The bounds are independent of polynomial degree, coefficient-field ramification and $\rho$ in the indicated range.
\end{theorem}

For example, at $p=2,h_0=1$, one member requires budget $H$, two require $\log_2 H+O(\log\log H)$, and three require $\log_2\log_2 H+O(\log\log\log H)$. In contrast, with no cardinality restriction the obstruction budget is $\sum_{h=h_0}^H p^{1-h}$, bounded as $H\to\infty$ (see below). These are statements about the order \emph{exponent} $h$; the torsion order is $p^h$. In particular the extra exponential converting exponents to orders must not be omitted.

\begin{theorem}[Exact finite optimization]\label{thm:finite}
For any $n,H$ and $\delta\ge0$, the infimum $C_n(H,\delta)$ is a minimum. It is independent of $\rho\ge r_{h_0}$ and is attained by monic polynomials all of whose roots belong to the tested shells, with thresholds $t_i=G_\rho(g_i)=v(g_i(0))$. No member needs more than $e_h$ roots, counting multiplicity, in shell $h$. There is an exact finite integer-state recurrence for the minimum, given in Section~\ref{sec:algorithm}. In particular,
\[
 C_n(H,\delta)=C_n(H,0)\qquad (0\le\delta<r_H).
\]
\end{theorem}

\section{Root accounting and the one-color endpoint}\label{sec:root}
\begin{lemma}[Exact root model]\label{lem:model}
For $g=c\prod_{\alpha}(T-\alpha)$ (roots with multiplicity),
\begin{align*}
 G_\rho(g)&=v(c)+\sum_\alpha\min(v(\alpha),\rho),\\
 v(g(z))-G_\rho(g)&=E_h(z)-\bigl(G_\rho(g)-G_{r_h}(g)\bigr),\quad z\in Z_h,\\
 E_h(z)&=\sum_{v(\alpha)=r_h}(v(z-\alpha)-r_h).
\end{align*}
If $m_{i,h}$ counts the shell-$h$ roots of member $i$, write $w_{i,h}=m_{i,h}r_h$, $W_i(h)=\sum_{k=h_0}^h w_{i,k}$. Then
\begin{equation}\label{eq:budget}
 \sum_i W_i(H)\le \sum_i t_i,\qquad
 G_\rho(g_i)-G_{r_h}(g_i)\ge S_i(h):=\sum_{k<h,\ k\ge h_0} w_{i,k}(1-p^{k-h}).
\end{equation}
\end{lemma}
\begin{proof}
The first formula is multiplicativity of the Gauss valuation (or the Gauss norm after scaling to the relevant radius). It holds at real radii by continuity in $r$ for polynomials. The ultrametric inequality is an equality for roots of valuation different from $v(z)$. This gives the second formula, including when $g(z)=0$. Integrality implies $G_0(g_i)\ge0$; hence $G_\rho(g_i)-G_0(g_i)$ includes the entire contribution $w_{i,h}$ from each tested shell. The contribution of a root in an earlier shell to $G_\rho-G_{r_h}$ is $r_k-r_h$. All other contributions are nonnegative. This proves \eqref{eq:budget}.
\end{proof}

For $\alpha$ with $v(\alpha)=r_h$, choose $\beta\in Z_h$ maximizing $v(\alpha-\beta)$. Then $v(z-\alpha)\le v(z-\beta)$ for every $z\in Z_h$: if $z\ne\beta$ and the reverse strict inequality held, the ultrametric equality would contradict maximality. Thus replacing on-shell roots by nearest torsion points can only increase on-shell excess. Other-shell excess depends only on their valuations. Removing all off-shell roots and the scalar $c$, making the new polynomial monic, cannot worsen loss: by Lemma~\ref{lem:model}, the discarded roots contributed only nonpositive amounts to $v(g(z))-G_\rho(g)$. Resetting thresholds to $G_\rho$ also cannot worsen loss and decreases budget. In the projected model, if any member has more than $e_h$ roots in shell $h$, replace those roots by the full shell, once each. This still gives infinite loss from that member throughout shell $h$; in all other shells the removed roots had nonpositive raw excess. Hence the modification preserves blockage and does not increase cost. The remaining finite possibilities prove the attainment and independence claims of Theorem~\ref{thm:finite}.

For completeness, the cyclotomic radius follows since all primitive $p^h$-th roots have $v(\zeta-1)$ equal (Galois invariance), and $\Phi_{p^h}(1)=p$. If $\xi$ is a primitive root and $a,b$ are distinct unit residues modulo $p^h$, then
\begin{equation}\label{eq:metric}
 v((\xi^a-1)-(\xi^b-1))=r_h p^{v_p(a-b)}.
\end{equation}
The nested unit congruence classes form a tree with $e_j$ classes at depth $j$.

\begin{lemma}[Shell capacity]\label{lem:capacity}
For $m<e_h$ let
\[
 U_h(m)=r_h\sum_{j=1}^{h-1}e_j\left\lfloor\frac m{e_j}\right\rfloor,
 \qquad U_h(e_h)=+\infty.
\]
The largest possible minimum of $E_h(z)$ over $Z_h$ for a choice of $m$ on-shell roots is $U_h(m)$, for $0\le m\le e_h$. It is attained by taking the first $m$ positive integers prime to $p$ as exponents of $\xi$.
\end{lemma}
\begin{proof}
Project roots to torsion points. Starting at the first congruence level, always choose a least-occupied class among the children of the current class. The chosen class at depth $j$ has at most $\lfloor m/e_j\rfloor$ marks, by induction. When $m<e_h$ the final leaf is unmarked. The contribution of a mark is the telescoping sum $p^d-1=\sum_{j=1}^d e_j$, where $d=v_p(a-b)$. Formula \eqref{eq:metric} now gives the upper bound. For the specified initial segment, each class at depth $j$ receives at least $\lfloor m/e_j\rfloor$ marks, giving the reverse inequality at every unmarked point; marked points have infinite excess. If $m=e_h$, use one root at every point.
\end{proof}

\begin{proposition}[One member, exact cost on a block]\label{prop:one}
Blocking consecutive shells $[a,b]$ at zero tolerance with one member costs at least $(b-a+1)p^{1-a}$, for any $a\ge h_0$. This is attained with exactly $(p-1)p^{h-a}$ roots in each shell $h\in[a,b]$ and no other roots.
\end{proposition}
\begin{proof}
Put $s=p^{1-a}$ and $M_h=(p-1)p^{h-a}$. Inductively suppose $m_k\ge M_k$ on all earlier tested shells. If $m_h<M_h$, the shell capacity is at most
\begin{align*}
 U_h(M_h-1)&=r_h\sum_{j=1}^{h-a}(M_h-e_j)\\
 &=(h-a)s-r_h(p^{h-a}-1).
\end{align*}
This includes $h=a$, when the sum is empty, and $a=1$. Indeed $M_h=e_{h-a+1}$, so all terms at depths $j\ge h-a+1$ vanish, and $e_j$ divides $M_h$ for smaller depths. The right hand side equals the contribution $\sum_{k=a}^{h-1}M_k(r_k-r_h)$ to the previous-root slack. Lemmas~\ref{lem:model}--\ref{lem:capacity} therefore give a good point, a contradiction. Induction shows $m_h\ge M_h$. Each shell costs $M_hr_h=s$.

Conversely mark the prescribed $M_h$ exponents on each shell and take the product of all linear factors. The minimum excess in shell $h$ is $U_h(M_h)=(h-a+1)s$ if $a>1$ and is infinite if $a=1$. This exceeds the earlier-root slack $(h-a)s-r_h(p^{h-a}-1)$. Later-shell roots contribute no slack. The monic product has budget exactly $(b-a+1)s$, and blocks the whole interval.
\end{proof}
This proves the $n=1$ formula of Theorem~\ref{thm:main}, independently of the article's equivalent endpoint. It also identifies the unlimited-cardinality cost at tolerance zero: a blocked shell needs at least $p-1$ roots collectively, since otherwise some first-level class is unmarked and has no excess for any polynomial (for $h=1$ the same assertion is the full-shell requirement). Using a separate optimal one-shell polynomial for each shell attains the lower bound $\sum_{h=h_0}^H p^{1-h}$.

\section{Delayed product potential: the upper bound}\label{sec:upper}
We first establish a robust bound on a consecutive interval after a given set of colors has become active. A color is active by level $a$ if it has at least one root in the tested shells through $a$. This is a root-count notion, not a condition on its values; an inactive polynomial can never have positive loss in a shell where it has no root. Write $C=B+1$ and assume $B\ge2$ for the bounds (replace a smaller budget by 2).

\begin{lemma}[A capped root average]\label{lem:cap}
For $v(\alpha)=r_h$,
\[
 \frac1{e_h}\sum_{z\in Z_h}\min\{v(z-\alpha)-r_h,C\}\le (h+C)r_h.
\]
\end{lemma}
\begin{proof}
Project to a nearest torsion point $\beta$. For the monic polynomial $P_h(T)=\Phi_{p^h}(1+T)$, the derivative at $\beta$ has valuation $h-r_1$: differentiate $(X^{p^h}-1)/(X^{p^{h-1}}-1)$ at $X=1+\beta$. Hence the sum of finite excesses at points other than $\beta$ is $h-r_1-(e_h-1)r_h\le h$. The term at $\beta$ is capped at $C$. Division by $e_h$ proves the bound. The calculation includes the one-point shell at $p=2,h=1$.
\end{proof}

\begin{lemma}[Forced growth on a blocked shell]\label{lem:forced}
If shell $h$ is blocked at zero tolerance and every color having roots in that shell has $W_i(h-1)>0$, then, summing over all colors active before $h$,
\begin{equation}\label{eq:forced}
 \sum_i\frac{w_{i,h}}{S_i(h)}>\frac1{h+C}.
\end{equation}
\end{lemma}
\begin{proof}
For each active $i$, $0<S_i(h)\le B$. At a blocked point $z$ at least one color has positive loss. By Lemma~\ref{lem:model} its $E_h(z)$ exceeds $S_i(h)$. Capping each root excess separately at $C>B$ preserves this strict inequality: either no term changes or a capped term alone is greater than $S_i(h)$. Thus the sum over active colors of the capped excess divided by $S_i(h)$ is strictly larger than one at every point. Take the uniform average and apply Lemma~\ref{lem:cap} to each root. Inactive colors with no root in the shell cannot have positive loss, so none are missing from this argument.
\end{proof}

\begin{lemma}[Delayed-potential interval bound]\label{lem:interval}
Suppose all shells from $a$ through $b$ are blocked and no new color is activated after $a$ through $b$. Let the $k$ colors active by $a$ have first activation levels $a_i\le a$. Let $d\ge1$ be an integer and $K>1$ a real number, and set
\[
 \lambda=\frac{1-p^{-d}}K,\quad L=\sum_{i=1}^k\log(B e_{a_i}),\quad
 D=d+\left\lceil\frac{dL}{\log K}\right\rceil.
\]
Then
\begin{equation}\label{eq:interval}
 b+C+1\le (a+C+1+D)\exp(L/\lambda).
\end{equation}
\end{lemma}
\begin{proof}
All logarithms in this proof are natural. For $h\ge a$ put $P_h=\prod_{i=1}^k W_i(h)>0$, and $R_h=P_h/P_{h-1}\ge1$ for $h>a$. Each $W_i(a)$ is at least $r_{a_i}$ and each $W_i(b)$ is at most $B$. Consequently $\sum_{h=a+1}^b\log R_h=\log(P_b/P_a)\le L$.

Discard the first $d$ shells after $a$ (or all if there are fewer). For the remaining $h$, call a shell irregular if $P_{h-1}/P_{h-d}>K$. The sum of $\log(P_{h-1}/P_{h-d})$ over these $h$ is at most $dL$, since every increment $\log R_j$ occurs at most $d$ times. There are therefore at most $\lceil dL/\log K\rceil$ irregular shells. On a regular shell, monotonicity of every $W_i$ and the bound on their product imply
\[
 W_i(h-d)\ge W_i(h-1)/K,
 \qquad S_i(h)\ge(1-p^{-d})W_i(h-d)\ge\lambda W_i(h-1).
\]
The middle inequality retains only the roots from levels at most $h-d$. Formula \eqref{eq:forced} now implies
\[
 R_h=\prod_i\left(1+\frac{w_{i,h}}{W_i(h-1)}\right)
 \ge 1+\sum_i\frac{w_{i,h}}{W_i(h-1)} > 1+\frac\lambda{h+C}.
\]
Since $0<\lambda<1$, concavity gives $\log(1+\lambda/x)\ge\lambda\log(1+1/x)$ for $x>0$. Summing over regular shells yields
\[
 L\ge\lambda\sum_{\text{regular }h}\log\left(1+\frac1{h+C}\right)
 \ge\lambda\log\frac{b+C+1}{a+C+1+D}.
\]
For the last inequality, the summands decrease with $h$, so removing at most $D$ terms removes no more than the first $D$ terms; if $D>b-a$, the right side is nonpositive and the inequality is automatic. Exponentiation gives \eqref{eq:interval}.
\end{proof}

The delayed comparison is essential to the leading constant: the elementary inequality $S_i(h)\ge(1-1/p)W_i(h-1)$ alone would put the factor $p/(p-1)$ in the exponent of the next activation bound. The exceptional-shell accounting allows the slack to approach $W_i$ instead.

\subsection*{Explicit recursive bound and asymptotic evaluation}
Consider any blocking packet for $[h_0,H]$, and let $m\le n$ be the number of colors having roots in the tested shells. Order their activation levels with multiplicity, $a_1\le\cdots\le a_m$; necessarily $a_1=h_0$. Adjoin $a_{m+1}=H+1$. If $m=1$, Proposition~\ref{prop:one} applies to the full interval. Otherwise, before $a_2$ only one color is active, and the proposition gives
\begin{equation}\label{eq:a2}
 a_2\le h_0+B/s_0,
\end{equation}
including the case of a tie $a_2=h_0$. On $[a_k,a_{k+1}-1]$ for $k\ge2$, precisely the first $k$ colors are active if the interval is nonempty; apply Lemma~\ref{lem:interval}. If the interval is empty the desired bound is automatic.

Here is an explicit recursion depending only on $p,h_0,n,B$ (use $B\gets\max(B,2)$ for finiteness at small budgets). Put
\[
 d=\lceil\log_p B\rceil,\quad K=1+B^{-1},\quad \lambda=(1-p^{-d})/K,
 \qquad U_1=h_0,\quad U_2=h_0+B/s_0.
\]
For $k=2,\ldots,n$ let
\begin{align*}
 L_k&= k\log(B(p-1))+\log(p)\sum_{i=1}^k(U_i-1),\\
 U_{k+1}&=(U_k+C+1+d+\lceil dL_k/\log K\rceil)\exp(L_k/\lambda)+1.
\end{align*}
The sequence is increasing, and induction using \eqref{eq:interval} and \eqref{eq:a2} gives $a_i\le U_i$ up to $a_{m+1}$, so every blocker satisfies $H+1\le U_{n+1}$ (with $U_2$ used for $n=1$). In particular the integer $\lceil U_{n+1}\rceil$ is a guaranteed endpoint. This proves existence at exact zero tolerance without using compactness or a tolerance limit.

As $B\to\infty$, $d=O_p(\log B)$, $\lambda^{-1}=1+O(B^{-1})$ and $(\log K)^{-1}=O(B)$. For $k=2$, the logarithm of the prefactor of $\exp(L_2/\lambda)$ is $O_{p,h_0}(\log B)$, while
\[
 L_2/\log p=(U_2-1)+(h_0-1)+2\log_p(B(p-1))
 =B/s_0+O_{p,h_0}(\log B).
\]
Hence $\log_p U_3=B/s_0+O_{p,h_0}(\log B)$. For fixed $k\ge3$, $L_k=O_{p,h_0,n}(U_k)$, $U_k\ge B$, and the logarithm of the prefactor is $O_{p,h_0,n}(\log U_k+\log B+\log(d+1))=O_{p,h_0,n}(\log U_k)$. Thus for some fixed $c_1\ge1$,
\[
 \log_p U_{k+1}\le c_1 U_k\qquad(k\ge3,\ B\ge2).
\]
Taking $k-2\ge1$ more logarithms changes the bound by at most an additive constant relative to taking $k-2$ logarithms of $U_k$ alone, for sufficiently large $B$. Iteration gives
\[
 \log_p^{\circ(n-1)} U_{n+1}\le B/s_0+O_{p,h_0,n}(\log B).
\]
The same bound applies to the least guaranteed endpoint $H_n(B)$; rounding has no effect on the asserted error. This proves the upper half of \eqref{eq:tower}.

\section{Matching lower bound and inverse cost}\label{sec:lower}
For $B$ sufficiently large choose $n$ consecutive blocks. Start the first at $a_1=h_0$, with length
\[
 \ell_1=\left\lfloor\frac{B-(n-1)}{s_0}\right\rfloor,\qquad a_2=h_0+\ell_1.
\]
For $2\le i\le n$, give block $i$ length $\ell_i=p^{a_i-1}$ and start the next block at $a_{i+1}=a_i+\ell_i$. Choose one polynomial per block as in Proposition~\ref{prop:one}. Each blocks its assigned shells regardless of the other polynomials' values, and the total budget is $s_0\ell_1+(n-1)\le B$. This packet blocks through $H=a_{n+1}-1$, so $H_n(B)\ge a_{n+1}$. Since
\[
 \log_p(a_i+p^{a_i-1})=a_i-1+o(1)\quad(a_i\to\infty),
\]
$n-1$ iterated logarithms of $a_{n+1}$ give $a_2+O_{p,n}(1)=B/s_0+O_{p,h_0,n}(1)$. This is the lower bound in \eqref{eq:tower}. It uses algebraic coefficients in one finite cyclotomic extension (which depends on $H$) and the common witness zero.

For the inverse cost, note that $C_n(H,0)\to\infty$ with $H$ by the already established finite bound at every fixed budget. Attainment gives the equivalence
\[
 C_n(H,0)\le B\quad\Longleftrightarrow\quad H_n(B)>H.
\]
The explicit construction yields $C_n(H,0)\le s_0 X_H+O(1)$: choose $B=s_0(X_H+c)$ with $c$ large enough that the lower bound in \eqref{eq:tower} guarantees $H_n(B)>H$. Substituting $B=C_n(H,0)$ into the upper bound of \eqref{eq:tower} and using $B=O(X_H)$ proves the other inequality of \eqref{eq:inverse}. Theorem~\ref{thm:main} follows.

\section{Exact finite costs, even at positive tolerance}\label{sec:algorithm}
We now complete Theorem~\ref{thm:finite} by spelling out a recurrence that is exact for arbitrary finite $H,n,\delta$. This is deliberately an exhaustive algorithm rather than a claim of a fast algorithm for astronomical endpoints. It improves on an optimization over unknown polynomials by reducing it to finite integer state spaces, with no dependence on coefficients or degree bounds. The accompanying \code{packet\_exact.py} implements it for small shells.

In the monic projected model, identify $Z_h$ with unit residues $a\pmod{p^h}$. In each color $i$ choose a multiset of $m_i\le e_h$ marks $b$, with repetition allowed. In units of $r_h$, its on-shell excess at leaf $a$ is
\[
 V_i(a)=\sum_{b\ \mathrm{marked}}(p^{v_p(a-b)}-1),
\]
interpreted as $+\infty$ when $a$ is a mark. For $m_i=e_h$ it suffices to use the entire shell, since it dominates every other arrangement of the same count. The state just after shell $h-1$ is the list of pairs $(M_i,K_i)$, where $M_i$ is the number of earlier roots and $K_i$ is their total valuation in units of $r_{h-1}$. At the first shell set both entries to zero. The prior slack at the new shell, in units of $r_h$, is exactly $pK_i-M_i$, since $r_{h-1}=p r_h$. No future root can contribute to this slack. Therefore a choice of layouts is feasible exactly when for every leaf $a$,
\begin{equation}\label{eq:dpcondition}
 \sum_i\pos{V_i(a)-(pK_i-M_i)}>\delta e_h.
\end{equation}
All finite terms on the left are nonnegative integers. Clip each term at $T_h=\lfloor\delta e_h\rfloor+1$ (use $T_h$ for infinity); then the condition is that every clipped sum be at least $T_h$. For every feasible layout transition to
\[
 (M_i',K_i')=(M_i+m_i,\ pK_i+m_i).
\]
The recurrence stores all reachable states and may identify permutations of the colors. The answer at $H$ is the minimum of $r_H\sum_i K_i$ over final states. Induction proves that the state description retains all information about past shells needed for the next shell. Every feasible path is realized by the product of its chosen torsion linear factors in each color, so the recurrence neither relaxes nor tightens the polynomial problem. The finite reduction of Section~\ref{sec:root} proves the converse. For rational $\delta$ this computation uses only integers and rationals; for general real $\delta$ only the finitely many integer cutoffs $T_h$ are required.

As all finite losses in the reduced model are integer multiples of $r_h$, blockage at zero tolerance in any shell $h\le H$ entails loss at least $r_h\ge r_H$ at every point, or infinite loss. Thus the same optimal packet works for any $\delta<r_H$. Monotonicity in tolerance proves the exact stabilization of $C_n(H,\delta)$ in Theorem~\ref{thm:finite}. It also proves the following consequence.
\begin{corollary}[No inverse-tolerance singularity for fixed cardinality]
For fixed $p,h_0,n,B$, $H_n(B,\delta)=H_n(B,0)$ for every $0\le\delta<r_{H_n(B,0)}$. In particular the least guaranteed exponent stays bounded as $\delta\downarrow0$.
\end{corollary}
\begin{proof}
Write $J=H_n(B,0)$, which is finite. If $J=h_0$ the claim is immediate. Otherwise a blocker of budget at most $B$ exists through $J-1$, and by attainment and stabilization a blocker also exists there at tolerance $\delta<r_J<r_{J-1}$. Thus $H_n(B,\delta)\ge J$. The reverse inequality follows from monotonicity.
\end{proof}

\subsection*{Independent small-case checks}
The code enumerates layouts with repeated roots, checks the capacity formula for 27 pairs of shell and root count, and solves 16 small instances of the packet recurrence. Some outputs are reproduced here (the archive includes the complete log). These computations are checks of the finite model, not evidence substituted for the asymptotic proof.
\begin{center}
\begin{tabular}{rrrrrrl}
\toprule
$p$ & $h_0$ & $H$ & $n$ & $\delta$ & $C_n$ & Case \\
\midrule
2&1&3&1&0&3&Single member\\
2&1&3&2&0&2&Intermediate\\
2&1&3&3&0&$7/4$&Unrestricted optimum\\
3&1&2&2&0&$4/3$&Odd prime\\
2&2&3&1&0&1&Later starting shell\\
2&2&3&2&0&$3/4$&Later starting shell\\
2&1&3&1&$1/4$&3&Positive tolerance\\
2&1&3&2&$1/4$&2&Positive tolerance\\
2&1&3&3&$1/4$&2&Positive tolerance\\
\bottomrule
\end{tabular}
\end{center}
The single-member cases reproduce Proposition~\ref{prop:one}; when $n\ge H-h_0+1$ the code reproduces the independent-shell costs from the article, which at positive tolerance are $\sum_h r_h\min\{m:U_h(m)>\delta\}$. The small zero-tolerance examples agree with the optimal consecutive-block partition upper bound. This does \emph{not} prove that every finite optimum is a consecutive-block partition; the recurrence includes mixed color assignments at individual shells.

\section{Dimension-free exact selection and a GRH application}\label{sec:multi}
\begin{theorem}[Common tuple at exact Gauss values]\label{thm:multi}
Let $F_1,\ldots,F_n$ be nonzero integral formal power series in $d\ge1$ variables over $\OO$ (polynomials included). Define $A_i=G_\rho(F_i)$ by the minimum of $v(c_\nu)+\rho|\nu|$ over coefficients. If $\sum_i A_i\le B$, then there exist $z_1,\ldots,z_d\in\bigcup_{h=h_0}^{H_n(B)} Z_h$ such that
\[
 v(F_i(z_1,\ldots,z_d))\le A_i\qquad(1\le i\le n).
\]
The least universal endpoint for this statement is exactly the one-variable endpoint $H_n(B)$, for every $d\ge1$. If instead witnesses $x_i\in\Cp^d$ satisfy $v((x_i)_j)\ge\rho$ and $\sum_i v(F_i(x_i))\le B<\infty$, the same endpoint works with $v(F_i(x_i))$ in place of $A_i$.
\end{theorem}
\begin{proof}
First consider polynomials. For each $i$, view $F_i$ as a polynomial in the remaining variables with coefficients $g_{i,\beta}(T_1)$. Choose $\beta_i$ for which $G_\rho(g_{i,\beta_i})+\rho|\beta_i|=A_i$; the coefficient polynomial is nonzero. All selected coefficient polynomials are integral and their total Gauss budget is at most $B$. The univariate theorem supplies a common $z_1$ with $v(g_{i,\beta_i}(z_1))\le G_\rho(g_{i,\beta_i})$ for every $i$. The specialized $F_i(z_1,\cdot)$ is nonzero and has Gauss valuation at most $A_i$ because this particular coefficient supplies an admissible upper bound. Its coefficients remain integral. Iterate through the other coordinates. In particular no loss accumulates with dimension, and the choice of $\beta_i$ may differ across members at every step.

For integral power series, $\rho>0$ implies the Gauss minimum is attained at finite total degree, since coefficient valuations are nonnegative. Truncate all series at a total degree $D$ so large that $D r_{H_n(B)}>\max_i A_i$. The truncations preserve $A_i$, and the omitted terms evaluate with valuation greater than $A_i$ at every tested tuple. Apply the polynomial argument and the ultrametric equality to transfer the inequalities back to the original series. Witness evaluations are convergent and satisfy $v(F_i(x_i))\ge A_i$, giving the last assertion. Sharpness follows by taking the univariate extremal packet, with no dependence on the other coordinates. For witness bounds the extremizers have witness $0$.
\end{proof}

The next consequence uses the supplied Dirichlet GRH, rather than just the local method. It is stated for ordinary Dirichlet characters; no unrestricted lifting of an arbitrary Selmer datum is asserted.
\begin{corollary}[Prime-conductor realization under Dirichlet GRH]\label{cor:characters}
Under the input Dirichlet GRH, for the hypotheses of Theorem~\ref{thm:multi} there exists a rational prime $q$ satisfying
\[
 q\equiv1\pmod{p^{H_n(B)}},\qquad q<2\bigl(p^{H_n(B)}H_n(B)\log p\bigr)^2,
\]
and a generator $u$ of $(\Z/q\Z)^\times$, and nonprincipal primitive Dirichlet characters $\chi_1,\ldots,\chi_d$ modulo $q$, valued in $\Cp$, of $p$-power order, for which
\[
 v\bigl(F_i(\chi_1(u)-1,\ldots,\chi_d(u)-1)\bigr)\le A_i\qquad(1\le i\le n).
\]
Each character has conductor $q$, its order lies between $p^{h_0}$ and $p^{H_n(B)}$, and the same statement holds with separate witnesses as in Theorem~\ref{thm:multi}.
\end{corollary}
\begin{proof}
Bach--Sorenson's explicit least-prime-in-progression bound under ERH \cite{bach} gives such a prime with modulus $p^{H_n(B)}$ and residue 1. This application concerns only Dirichlet $L$-functions (the cyclotomic field is abelian). The unit group of a prime field is cyclic. For the roots $1+z_j$ supplied by Theorem~\ref{thm:multi}, define the character by $\chi_j(u)=1+z_j$. The divisibility $p^{H_n(B)}\mid q-1$ ensures each defines a homomorphism on the entire cyclic group. None is principal since $h_0\ge1$; thus the conductor of each is the prime $q$. Substitution proves the result. Identifying the finite character values with complex roots of unity if necessary does not change which Dirichlet GRH input is used in the prime bound.
\end{proof}
The conductor bound is coarse and not claimed sharp. The underlying local exponent is sharp. This provides a concrete arithmetic realization of the zero-tolerance specialization, but does not ensure other prescribed Frobenius, ramification, or Selmer constraints in the original proof.

\section{Novelty, scope and further research}
\textbf{What has been established.} The main theorem identifies the sharp leading coefficient and the exact number of iterated logarithms in the fixed-cardinality obstruction cost. It answers the zero-tolerance asymptotic part of the explicit open question \code{cost:q:cardinality}, also giving a complete finite (possibly expensive) procedure for all tolerances. The zero-tolerance packet selector works in arbitrary dimension without multiplying the loss by the number of coordinates. The argument does not require the truth of RH; the character-conductor corollary explicitly does use Dirichlet GRH. A focused literature search for simultaneous torsion polynomial valuation bounds and fixed-cardinality optimal levels located the source-related workbook and related nonvanishing literature, but did not locate the tower law proved here. This is evidence for novelty relative to the inspected sources, not a comprehensive priority claim.

\textbf{Limitations.} We did not independently replay the 818-file vendored Lean package, audit all 914 article statements, or compile a Lean formalization of this new theorem. The source README's reported checks are distinct from our exact rational tests. The proof concerns unrestricted coefficient fields: the lower examples can require highly ramified cyclotomic coefficients. The exact finite computation enumerates multisets and is exponential or worse as shell sizes grow. No conclusion is claimed about the existence of all arithmetic twists required for the upstream Selmer argument, or about unresolved nonabelian GRH statements.

\textbf{Further research questions.}
\begin{enumerate}
\item Can the $O(\log B)$ upper error in \eqref{eq:tower} be reduced to $O(1)$, or can the next term be determined for $n=2$? Is the optimum always obtainable by consecutive blocks, or do mixed color layouts beat it?
\item At fixed positive $\delta$, determine the exact asymptotic of $H_n(B,\delta)$ as $B\to\infty$, and the joint transition as $\delta$ decreases with $B$. The finite recurrence is a way to test candidate formulas; the zero-tolerance tower law alone does not give the positive-tolerance asymptotic.
\item Impose Galois stability of each root multiset over a fixed finite coefficient field, as in the article's field-specific costs. Determine how ramification and cyclotomic intersections constrain activation and alter the tower height.
\item Extend the delayed potential to formal groups of finite height and to ramified Lubin--Tate settings, tracking shell radii and overlapping valuations. The cyclotomic proof uses equal spacing ratios and precise finite-shell geometry.
\item Formalize the root model, shell capacity, finite recurrence and delayed product potential in Lean. The existing \code{RootsOfUnity.lean} is a natural entry point. Then test whether any fixed-cardinality specialization can be inserted into a particular arithmetic Selmer twisting step with the necessary local and global constraints.
\end{enumerate}

\begin{thebibliography}{99}
\bibitem{vendor} ProveIt, \emph{Vendored grh-lean (readable edition)}, \href{https://github.com/VladimirReshetnikov/ProveIt/tree/main/Analysis/GeneralizedRiemannHypothesis/vendor/grh-lean}{repository directory}; vendor README blob \texttt{a04f66a62af93fc0757f2aa573f3b053fde51bfa}.
\bibitem{leanreadme} L. Sernau's repository, \emph{The Generalized Riemann Hypothesis for Dirichlet L-functions, in Lean 4}, readable edition, \href{https://github.com/lsernau/grh-lean/tree/72fead102f663d5c7a407964a13dc0cf9d977207/grh}{pinned upstream commit}. Proof credited by the repository to L. Alp\"oge with Claude.
\bibitem{statements} Vendored \href{https://github.com/VladimirReshetnikov/ProveIt/blob/main/Analysis/GeneralizedRiemannHypothesis/vendor/grh-lean/grh/src/Statements.lean}{\texttt{src/Statements.lean}}, blob \texttt{1cd1bce93344065936b2b392b23306254e0785ef}.
\bibitem{main} Vendored \href{https://github.com/VladimirReshetnikov/ProveIt/blob/main/Analysis/GeneralizedRiemannHypothesis/vendor/grh-lean/grh/src/Main.lean}{\texttt{src/Main.lean}}, blob \texttt{1adf4187a938510363cdb0dde0a6d628ff39c600}.
\bibitem{inequality} Vendored \href{https://github.com/VladimirReshetnikov/ProveIt/blob/main/Analysis/GeneralizedRiemannHypothesis/vendor/grh-lean/grh/src/New/Reduction/MainInequality.lean}{\texttt{MainInequality.lean}}, blob \texttt{97f83ab52b2feb420e952a2722b4cef454b9e579}.
\bibitem{onesided} Vendored \href{https://github.com/VladimirReshetnikov/ProveIt/blob/main/Analysis/GeneralizedRiemannHypothesis/vendor/grh-lean/grh/src/New/Reduction/OneSidedBound.lean}{\texttt{OneSidedBound.lean}}, blob \texttt{99e5240de112527b739068b692e5007720d310c1}.
\bibitem{roots} Vendored \href{https://github.com/VladimirReshetnikov/ProveIt/blob/main/Analysis/GeneralizedRiemannHypothesis/vendor/grh-lean/grh/src/New/Congruences/Selmer/RootsOfUnity.lean}{\texttt{RootsOfUnity.lean}}, blob \texttt{b5ad128a2c44538e8783050a371e1e44af0c32d1}.
\bibitem{twisting} Vendored \href{https://github.com/VladimirReshetnikov/ProveIt/blob/main/Analysis/GeneralizedRiemannHypothesis/vendor/grh-lean/grh/src/New/Congruences/Selmer/TwistingTheorem.lean}{\texttt{TwistingTheorem.lean}}, blob \texttt{3f5a5fc163a22758b93731ac87cc333e2f981b70}.
\bibitem{article} ProveIt, \emph{Beyond the Riemann hypothesis for Dirichlet L-functions: consequences, transfers and limits}, \href{https://github.com/VladimirReshetnikov/ProveIt/tree/main/Analysis/GeneralizedRiemannHypothesis/Article}{merged article}, README blob \texttt{7f3286088c94a8df3d8b9846112ea83dbbbbfa25}; see especially \code{sections/src.tex}, blob \texttt{f69668b2f95261630bb1282c2cc1f9c172b7e530}.
\bibitem{artcost} Ibid., \href{https://github.com/VladimirReshetnikov/ProveIt/blob/main/Analysis/GeneralizedRiemannHypothesis/Article/sections/cost.tex}{\texttt{sections/cost.tex}}; see labels \code{cost:thm:minimax}, \code{cost:thm:packet}, \code{cost:thm:transition}, \code{cost:q:cardinality}.
\bibitem{arttors} Ibid., \href{https://github.com/VladimirReshetnikov/ProveIt/blob/main/Analysis/GeneralizedRiemannHypothesis/Article/sections/tors.tex}{\texttt{sections/tors.tex}}; see labels \code{tors:thm:selection}, \code{tors:thm:fraction}, \code{tors:thm:packet}, \code{tors:q:cardinality}.
\bibitem{workbook} \emph{Roots of unity for many polynomials}, GRH workbook, 26 November 2026, \href{https://ultronozm.github.io/grh-workbook/20261126T040300--roots-of-unity-for-many-polynomials__workbook.html}{online chapter} (accessed 27 November 2026).
\bibitem{bach} E. Bach and J. Sorenson, Explicit bounds for primes in residue classes, \emph{Math. Comp.} \textbf{65} (1996), 1717--1735, \href{https://doi.org/10.1090/S0025-5718-96-00763-6}{doi:10.1090/S0025-5718-96-00763-6}.
\end{thebibliography}
\end{document}
